\section{Effective analysis: from local certificates to coherent names}
\label{sec:analysis}

Many analysis packets correctly replace limits by rational approximants and
error bounds, but a repeated gap remains: separately certified approximants
are treated as one coherent real name.  FIBER distinguishes three levels:
pointwise approximation, finite-window stability, and a global Cauchy modulus.
Only the third is a completed name.

\subsection{Dyadic names and adjacent motion}

Represent a dyadic rational at precision \(n\) by a signed integer
\(a_n\), interpreted as \(a_n/2^n\).  Signed integers can be represented by the
existing difference terms in \texttt{IPRADerivations.Core.Integers}.  Define
the scaled discrepancy, for \(m\ge n\), by
\[
       \Delta(n,m)=\left|2^{m-n}a_n-a_m\right|.
\]
This is an integer quantity and therefore avoids division in the object
language.

\begin{theorem}[Dyadic adjacent-motion theorem]
\label{thm:dyadic}
Suppose a sequence of signed integers satisfies
\[
        |2a_k-a_{k+1}|\le 1
        \qquad\text{for every }k.
\]
Then for all \(m\ge n\),
\[
        \Delta(n,m)\le 2^{m-n}-1.
\]
Consequently
\[
        \left|\frac{a_n}{2^n}-\frac{a_m}{2^m}\right|<2^{-n}.
\]
Thus \(n\mapsto a_n/2^n\) is a Cauchy name with modulus \(n\mapsto n\).
\end{theorem}

\begin{proof}
Write \(r=m-n\) and induct on \(r\).  For \(r=0\), the discrepancy is zero,
which equals \(2^0-1\).  For the successor step, use
\begin{align*}
\Delta(n,m+1)
 &=|2^{m+1-n}a_n-a_{m+1}|\\
 &=|2(2^{m-n}a_n-a_m)+(2a_m-a_{m+1})|\\
 &\le 2\Delta(n,m)+|2a_m-a_{m+1}|\\
 &\le 2(2^{m-n}-1)+1\\
 &=2^{m+1-n}-1.
\end{align*}
The induction invariant is an inequality between integer-coded terms, hence
quantifier-free.  The object proof uses integer triangle inequality, powers of
two, order monotonicity, and quantifier-free induction.

Dividing the scaled inequality by \(2^m\) is only an explanatory
interpretation.  The formal endpoint can remain
\(\Delta(n,m)<2^{m-n}\), an integer inequality.  This is exactly the usual
Cauchy estimate after clearing denominators.
\end{proof}

\begin{corollary}[Independent certificates become coherent]
Suppose that for each precision \(n\) a local certificate produces \(a_n\), and
a pair certificate proves \(|2a_n-a_{n+1}|\le 1\).  Then the certificates may
be sealed while the sequence of observations exports a Cauchy name.
\end{corollary}

\begin{proof}
Use local existential elimination to obtain each adjacent pair only inside the
proof of its inequality.  The public conclusion depends solely on the
observations \(a_n,a_{n+1}\).  Apply \cref{thm:dyadic}.  No single global
certificate choosing all pointwise witnesses is necessary unless the name
itself must be returned as a code; in that case bounded least selection and a
master certificate provide one.
\end{proof}

\subsection{General summable-tail theorem}

The dyadic estimate is a special case of a finite telescoping principle.
Let \(e_k\ge0\) be integer or rational error majorants, and let \(T(n)\) be a
primitive-recursive tail certificate satisfying, for every \(m\ge n\),
\[
        \sum_{k=n}^{m-1}e_k\le T(n).
\]
If \(|x_{k+1}-x_k|\le e_k\), then
\[
        |x_m-x_n|\le T(n).
\]
The proof is finite induction using the triangle inequality and the bounded
sum equations.  A global rate therefore follows from an explicit summable
tail majorant; it does not follow from pointwise convergence or metastability
alone.

\subsection{Metastability is a separate interface}

A metastable contract has the form
\[
 \forall k,g\;\exists N\le\Phi(k,g)\;
 \forall i,j\in[N,N+g(N)]\;d(x_i,x_j)\le2^{-k}.
\]
For fixed \(k,g\) this is a finite query and fits FIBER directly: a certificate
contains \(N\) and a finite table or checker for the window.  It is often
provable when no useful global rate is known.  The framework keeps it as a
metastable interface.  Converting it to a Cauchy name requires additional
data, such as a summable adjacent-motion estimate or an independently supplied
tail modulus.  This prevents an invalid strengthening found in several
informal proof plans.

\subsection{Finite-net lifting}

Abstract compactness should be replaced by the finite data actually consumed
at a target precision.

\begin{theorem}[Finite-net lifting]
\label{thm:net}
Fix encoded domains \(X,Y\), discrepancy functions \(d_X,d_Y\), a function
\(f:X\to Y\), and a predicate \(P(y)\).  For a target query suppose there are:

\begin{enumerate}
\item a finite list \(z_0,\ldots,z_{L-1}\) and a coverage certificate
\[
       \forall x\in X\;\exists i<L\;d_X(x,z_i)\le\epsilon;
\]
\item a modulus certificate
\[
       d_X(x,z)\le\epsilon\longrightarrow
       d_Y(f(x),f(z))\le\delta;
\]
\item a stability theorem
\[
       P(f(z))\wedge d_Y(f(x),f(z))\le\delta
       \longrightarrow P_\delta(f(x));
\]
\item finite proofs of \(P(f(z_i))\) for all \(i<L\).
\end{enumerate}

Then \(\ipra\) proves \(P_\delta(f(x))\) for every encoded \(x\in X\).
The chosen net point is local and need not be exported.
\end{theorem}

\begin{proof}
Introduce an arbitrary \(x\) satisfying the domain predicate.  Open the
coverage existential to obtain an index \(i<L\) with the distance bound.
Instantiate the finite net theorem at \(i\), apply the modulus, then apply
stability.  The conclusion does not mention \(i\), so existential elimination
seals it.  The finite family of net-point proofs can be represented by a
master certificate or by a reflected bounded universal.  Every induction used
to verify list lookup or the finite family is quantifier-free after
Booleanization.
\end{proof}

\begin{remark}
This theorem does not assert that an arbitrary compact metric space has an
effective net in the selected coding.  Net production and coverage are
explicit hypotheses or residuals.  Once supplied, the remainder of the usual
compactness argument is a finite object-theory proof.
\end{remark}

\subsection{Guarded partial operations}

A statement such as \(d\ne0\) is generally too weak for stable numerical
division.  The useful interface is a positive apartness margin:
\[
             |d|\ge\eta>0.
\]
At target precision \(n\), the demand transformer requests enough precision on
\(d\) and the numerator to keep interval error below a fixed fraction of
\(\eta\).  Interval or Newton certificates then verify the reciprocal and the
quotient.  The same pattern applies to pivots, logarithms, square roots away
from zero, normalized eigenvectors, and branch labels.  The guard is a
residual if no effective lower bound is known.

\subsection{Modulus calculus}

FIBER realizers carry a demand map, which functions as a compositional modulus.
For metric interfaces, common operations have the following finite rules.

\begin{center}
\begin{tabularx}{0.96\linewidth}{Y{0.29\linewidth}X}
\toprule
Operation & Sufficient upstream demand\\
\midrule
addition & enough precision that the two errors sum below the target\\
multiplication on a bounded box & input errors divided by a certified magnitude bound\\
composition \(g\circ f\) & first apply the modulus of \(g\), then that of \(f\)\\
finite sum of \(L\) terms & target precision plus a binary-length allowance for \(L\)\\
division & target precision plus an apartness margin for the denominator\\
root label & cluster precision plus a positive separation gap\\
finite integral or quadrature & mesh and local error sufficient for the requested total error\\
\bottomrule
\end{tabularx}
\end{center}

These are not semantic continuity assumptions hidden in a typeclass.  Each is
a finite arithmetic contract with an object derivation.  Composition of the
demand maps is the proof-mining content of the framework.

\subsection{Worked analytic pattern: certified bisection}

For a rational polynomial \(p\), an interval \([l_0,r_0]\), and a target
precision \(n\), a certificate contains the bisection trace
\((l_k,r_k)_{k\le n}\).  The verifier checks:
\[
 p(l_k)\le0\le p(r_k),\qquad
 r_k-l_k=2^{-k}(r_0-l_0),
\]
and verifies that \((l_{k+1},r_{k+1})\) is the half selected by the exact
rational sign of the midpoint.  These are finite polynomial and order
calculations.  A uniqueness claim additionally requires a guard: for example,
a certificate that
\[
 p(y)-p(x)=(y-x)Q(x,y)
 \quad\text{and}\quad Q(x,y)\ge\lambda>0
\]
throughout the interval.  Then two zeros satisfy
\((y-x)Q(x,y)=0\), and positive cancellation gives \(x=y\).  This avoids an
appeal to the classical mean-value theorem and fits the checked ring/order
library.
